\documentclass[10pt,twocolumn,letterpaper]{article}
\usepackage[T1]{fontenc}

\usepackage{amsmath,amssymb,amsthm,mathtools,bm,braket}
\usepackage{mathptmx}
\renewcommand{\ttdefault}{pcr}
\usepackage[margin=0.68in,columnsep=0.26in,headheight=14pt,headsep=0.15in,footskip=0.28in]{geometry}
\usepackage{graphicx,booktabs,array,tabularx,longtable,multirow}
\usepackage{microtype,fancyhdr,caption,titlesec,enumitem,float,placeins}

\usepackage[hidelinks,bookmarksnumbered=true,hypertexnames=false]{hyperref}
\usepackage{xurl}
\setlength{\parindent}{1em}
\setlength{\parskip}{0.16em}
\setlength{\emergencystretch}{2em}
\setlength{\abovedisplayskip}{6pt plus2pt minus2pt}
\setlength{\belowdisplayskip}{6pt plus2pt minus2pt}
\captionsetup{font=small,labelfont=bf,justification=raggedright,singlelinecheck=false}
\titleformat{\section}{\large\bfseries}{\thesection.}{0.45em}{}
\titleformat{\subsection}{\normalsize\bfseries}{\thesubsection}{0.45em}{}
\titleformat{\subsubsection}{\normalsize\bfseries}{\thesubsubsection}{0.45em}{}
\titlespacing*{\section}{0pt}{11pt plus2pt minus2pt}{5pt}
\titlespacing*{\subsection}{0pt}{8pt plus2pt minus2pt}{3pt}
\pagestyle{fancy}\fancyhf{}
\fancyhead[L]{\footnotesize SAM Volume \VolumeRoman\ | \leftmark}
\fancyhead[R]{\footnotesize Brady | 19 October 2026}
\fancyfoot[C]{\footnotesize\thepage}
\renewcommand{\headrulewidth}{0.3pt}
\newtheorem{theorem}{Theorem}[section]
\newtheorem{proposition}{Proposition}
\newtheorem{lemma}[theorem]{Lemma}
\newtheorem{corollary}[theorem]{Corollary}
\theoremstyle{definition}
\newtheorem{definition}[theorem]{Definition}
\newtheorem{assumption}[theorem]{Assumption}
\newtheorem{example}[theorem]{Example}
\theoremstyle{remark}
\newtheorem{remark}[theorem]{Remark}
\newcommand{\dd}{\mathrm{d}}
\newcommand{\ii}{\mathrm{i}}
\newcommand{\Tr}{\operatorname{Tr}}
\newcommand{\Span}{\operatorname{span}}
\allowdisplaybreaks[2]
\raggedbottom
\usepackage{tocloft}
\setlength{\cftsecnumwidth}{3.2em}



\setcounter{tocdepth}{1}
\renewcommand{\thesection}{\Roman{section}}
\renewcommand{\thesubsection}{\thesection.\arabic{subsection}}
\renewcommand{\thesubsubsection}{\thesubsection.\arabic{subsubsection}}
\numberwithin{equation}{section}
\numberwithin{figure}{section}
\numberwithin{table}{section}
\newcommand{\bookchapter}[1]{%
\clearpage\refstepcounter{section}\setcounter{subsection}{0}%
\addcontentsline{toc}{section}{\protect\numberline{\thesection}#1}%
\markboth{Chapter \thesection}{}%
\twocolumn[{\begin{center}{\small SAM VOLUME \VolumeRoman\par}\vspace{0.5em}
{\LARGE\bfseries #1\par}\vspace{0.4em}{\large Chapter \thesection\par}
\vspace{0.65em}\end{center}}]}
\providecommand{\norm}[1]{\left\lVert#1\right\rVert}
\providecommand{\rank}{\operatorname{rank}}
\providecommand{\sgn}{\operatorname{sgn}}
\providecommand{\diag}{\operatorname{diag}}
\providecommand{\Var}{\operatorname{Var}}
\providecommand{\Cov}{\operatorname{Cov}}
\providecommand{\tightlist}{\setlength{\itemsep}{0pt}\setlength{\parskip}{0pt}}
\providecommand{\passthrough}[1]{#1}
\providecommand{\pandocbounded}[1]{#1}
\providecommand{\real}{\mathbb R}
\providecommand{\Complex}{\mathbb C}

\newcommand{\samref}[1]{\cite{#1}}

\newcommand{\bookpart}[2]{%
\clearpage\onecolumn\markboth{Part #1}{}%
\phantomsection\addcontentsline{toc}{part}{Part #1: #2}%
\vspace*{\fill}\begin{center}{\Large\bfseries Part #1\par}\vspace{0.4in}
{\LARGE\bfseries #2\par}\end{center}\vfill\clearpage}
